% \section{Optimization-Based Prompt Recovery Attack}
% \label{sec:attacj}
% To systematically evaluate whether erased concepts can be recovered without explicit lexical references, we design an automated prompt search procedure that identifies natural-language prompts capable of reactivating the target concept. Our attack operates under the threat model defined in Section~\ref{sec:threat_model}, where the adversary has query access to the unlearned diffusion model but is prohibited from explicitly using the target token.

% \subsection{Attack Objective}

% Let $c$ denote a target concept (e.g., \textit{horse}) and $t_c$ the corresponding token removed during unlearning. Given an unlearned diffusion model $D_{\theta'}$, the adversary seeks a prompt $p$ that satisfies two conditions:

% \begin{itemize}
% \item the prompt does not contain the target token $t_c$
% \item the generated image contains the concept $c$
% \end{itemize}

% Formally, the adversary aims to solve

% \begin{equation}
% p^* = \arg\max_{p \in \mathcal{P},, t_c \notin p}
% \mathbb{E}*{x \sim p*{\theta'}(\cdot|p)}[S_c(x)],
% \end{equation}

% where $S_c(x)$ is a concept presence score measuring whether the generated image $x$ contains concept $c$. In practice, this score can be implemented using an external detector such as a CLIP similarity score or a trained concept classifier.

% \subsection{Prompt Search Space}

% Direct optimization over arbitrary text strings often leads to unnatural prompts. To ensure realistic attack scenarios, we restrict the search space to natural-language prompts derived from contextual tokens associated with the target concept.

% Specifically, we construct a candidate token set $S_c$ consisting of words that frequently co-occur with concept $c$ in training captions or semantic neighborhoods. For example, for the concept \textit{horse}, this set may include tokens such as \textit{cowboy}, \textit{ranch}, \textit{stable}, \textit{rider}, and \textit{fence}. Prompts are then generated by composing subsets of tokens from $S_c$ with natural-language templates.

% This design reflects the hypothesis that diffusion models encode concepts through distributed contextual representations rather than isolated lexical triggers.

% \subsection{Prompt Optimization}

% Starting from an initial set of seed prompts, the attack iteratively refines prompts to maximize the concept score. At each iteration, candidate prompts are generated by editing existing prompts through token substitutions, insertions, or deletions within the contextual token set $S_c$. Each candidate prompt is evaluated by generating images from the unlearned model and computing the concept score.

% The highest-scoring prompts are retained and used to generate new candidates in the next iteration. This process continues until convergence or a predefined search budget is exhausted.

% \begin{algorithm}[t]
% \caption{Natural Prompt Recovery Attack}
% \begin{algorithmic}[1]
% \STATE Initialize seed prompt set $\mathcal{P}*0$
% \FOR{$k = 1 \dots K$}
% \STATE Generate candidate prompts by editing prompts in $\mathcal{P}*{k-1}$
% \STATE Generate images using $D_{\theta'}$ for each prompt
% \STATE Compute concept score $S_c(x)$ for generated images
% \STATE Select top-scoring prompts to form $\mathcal{P}_k$
% \ENDFOR
% \STATE Return highest-scoring prompt $p^*$
% \end{algorithmic}
% \end{algorithm}

% \subsection{Implementation Details}

% In our implementation, the concept score $S_c(x)$ is computed using a CLIP-based similarity between the generated image and a textual description of the target concept. Prompts containing the banned token $t_c$ are filtered during candidate generation. To maintain prompt fluency, edits are restricted to semantically related tokens and common prompt templates.

% This optimization procedure allows the adversary to discover prompts that indirectly activate the erased concept, revealing weaknesses in token-level unlearning objectives.

\section{Optimization-Based Prompt Recovery Attack}
\label{sec:attack}

We study whether concept unlearning removes a concept itself or merely blocks \emph{direct lexical access} to that concept. To this end, we introduce an optimization-based prompt recovery attack that searches for \emph{indirect natural-language prompts} capable of recovering the erased concept without explicitly mentioning it.

\subsection{Threat Model}
\label{subsec:threat_model}

Let $G_u$ denote a concept-unlearned text-to-image diffusion model, where a target concept $c$ (e.g., \texttt{horse}) is intended to be forgotten. We consider a black-box attacker who can query $G_u$ with text prompts and observe the generated images, but cannot access model parameters or gradients. The attacker is not allowed to explicitly mention the erased concept in the prompt. The attacker's goal is to construct a prompt $p$ such that the generated image $I = G_u(p)$ still contains the target concept $c$.

This setting models a realistic failure mode of concept unlearning: a model may appear to forget a concept under standard direct-prompt evaluation, yet still recover that concept when prompted indirectly through semantically related scenes, roles, actions, or co-occurring objects.

\subsection{Constrained Prompt Optimization Objective}
\label{subsec:attack_objective}

We formulate indirect concept recovery as a constrained black-box optimization problem over natural-language prompts. Given a prompt search space $\mathcal{P}$, we seek a prompt
\begin{equation}
p^* = \arg\max_{p \in \mathcal{P}} \; \mathcal{J}(p; c),
\end{equation}
where the objective $\mathcal{J}$ encourages recovery of the target concept while discouraging trivial lexical leakage and unnatural prompts:
\begin{equation}
\mathcal{J}(p; c)
=
S_{\mathrm{target}}(G_u(p), c)
-
\lambda \, S_{\mathrm{explicit}}(p, c)
+
\gamma \, S_{\mathrm{natural}}(p).
\label{eq:attack_objective}
\end{equation}

Here, $S_{\mathrm{target}}$ measures the extent to which the generated image contains the target concept $c$, $S_{\mathrm{explicit}}$ penalizes prompts that explicitly mention or too directly reveal the erased concept, and $S_{\mathrm{natural}}$ encourages fluent, human-plausible prompts. Scalars $\lambda,\gamma \ge 0$ control the tradeoff between these terms.

\paragraph{Target concept recovery score.}
Given a generated image $I = G_u(p)$, we define a target recovery score $S_{\mathrm{target}}(I,c)$ to quantify how strongly $I$ expresses concept $c$. In practice, this score can be instantiated using one or more surrogate judges, such as:
\begin{equation}
S_{\mathrm{target}}(I,c)
=
\alpha \, S_{\mathrm{CLIP}}(I,c)
+
\beta \, S_{\mathrm{VLM}}(I,c)
+
\delta \, S_{\mathrm{det}}(I,c),
\label{eq:target_score}
\end{equation}
where $S_{\mathrm{CLIP}}$ denotes image-text similarity to concept descriptions of $c$, $S_{\mathrm{VLM}}$ denotes a vision-language model judge score, and $S_{\mathrm{det}}$ denotes an object detector or classifier confidence when available.

\paragraph{Explicit mention penalty.}
To prevent trivial attacks that directly name the forgotten concept, we impose a lexical penalty:
\begin{equation}
S_{\mathrm{explicit}}(p,c)
=
\mathbb{I}[\text{$p$ contains the exact target token}]
+
\eta \, \mathrm{LexSim}(p,c),
\label{eq:explicit_penalty}
\end{equation}
where the indicator term excludes prompts that explicitly mention $c$ (e.g., \texttt{horse}), and $\mathrm{LexSim}(p,c)$ optionally penalizes prompts that are lexically too close to the target concept. In our primary setting, we enforce an exact-token ban and treat stricter lexical constraints as an ablation.

\paragraph{Prompt fluency.}
We include a fluency score to favor interpretable prompts that remain close to realistic user inputs:
\begin{equation}
S_{\mathrm{fluent}}(p) = \mathrm{LMScore}(p),
\end{equation}
where $\mathrm{LMScore}$ may be instantiated by language-model likelihood or a fluency judge. 

\subsection{Prompt Search Space}
\label{subsec:prompt_space}

A key challenge is that stronger unlearning methods may resist shallow bag-of-words triggers while remaining vulnerable to richer semantic compositions. We therefore search over a structured space of indirect prompts that capture multiple forms of semantic association with the target concept:
\begin{itemize}
    \item \textbf{Context prompts}: scenes strongly associated with the target concept (e.g., ranch, stable, pasture).
    \item \textbf{Role/action prompts}: professions, agents, or activities commonly involving the target concept (e.g., rider, cowboy, carriage scene).
    \item \textbf{Co-occurrence prompts}: objects or attributes frequently appearing with the target concept (e.g., saddle, reins, hay bale).
    \item \textbf{Compositional prompts}: multi-cue prompts that combine scene, role, action, and object information in a natural description.
\end{itemize}

This richer search space is important for probing robustly unlearned models, which may suppress direct lexical triggers but still retain latent access paths through correlated semantic cues.

\subsection{Search Algorithm}
\label{subsec:search_algorithm}

Because the attacker does not observe gradients, we optimize Eq.~\ref{eq:attack_objective} using black-box prompt search. Starting from an initial pool of seed prompts, the search iteratively proposes modified prompts, queries the unlearned model, scores the generated outputs, and retains the most promising candidates.

Let $\mathcal{B}_t$ denote the candidate pool at iteration $t$. At each step, we generate a set of edited prompts $\widetilde{\mathcal{B}}_t$ through local transformations such as token substitution, phrase insertion, scene replacement, and action-object recombination. We then evaluate all candidates using Eq.~\ref{eq:attack_objective} and keep the top-$k$ prompts:
\begin{equation}
\mathcal{B}_{t+1} = \mathrm{TopK}\bigl(\mathcal{B}_t \cup \widetilde{\mathcal{B}}_t,\; \mathcal{J}(\cdot;c),\; k\bigr).
\end{equation}
After $T$ rounds, the final attack prompt is selected as the highest-scoring candidate in $\mathcal{B}_T$.

In practice, this search can be implemented using beam search, evolutionary search, or other query-based optimizers. Our implementation uses $\langle \textit{replace with beam search / evolutionary search} \rangle$.
